
\chapter{Introdução}

Consolidada como um produto extremamente popular, a televisão é o dispositivo central na maioria das salas de estar da população mundial. Em decorrência disso, a população se adaptou à prática de consumo de conteúdo em vídeo, criando uma base de consumidores preparada para a revolução do entretenimento digital. Com a popularização da Internet e o surgimento das plataformas de \textit{streaming}, esse hábito de consumo foi potencializado, migrando para um modelo de consumo sob demanda. Essa mudança gerou um crescimento exponencial no volume de dados de vídeo trafegando pela rede, impondo uma pressão sem precedentes sobre a infraestrutura global e tornando o balanceamento de tráfego uma necessidade crítica para garantir a qualidade do serviço aos usuários finais \cite{cisco2020vni}.

Desse modo, pode-se constatar que há desafios no contexto da otimização contínua dos processos de distribuição de conteúdo sob demanda, haja vista que é necessário maximizar a Qualidade de Experiência (QoE) oferecida pelos serviços para cada vez mais usuários. Contudo, é crucial reconhecer que a percepção de qualidade de experiência pelo usuário final (QoE) não depende unicamente da infraestrutura e capacidade dos provedores de serviço. Ela é, na verdade, uma responsabilidade compartilhada, na qual os provedores de acesso à Internet (ISPs) desempenham um papel igualmente fundamental. Frequentemente, a infraestrutura de rede dos ISPs não acompanha a evolução dos provedores de serviço e, por isso, a qualidade percebida pelo usuário pode não ser ótima. 

A distribuição de conteúdo de vídeo em larga escala geralmente ocorre por meio de uma arquitetura complexa e geograficamente dispersa, conhecida como Rede de Distribuição de Conteúdo (CDN - \textit{Content Delivery Network}). Nessa estrutura, o conteúdo original é replicado do servidor de origem para múltiplos servidores de borda (\textit{edge servers}) localizados em pontos estratégicos, mais próximos dos usuários finais. Quando um espectador solicita um vídeo, a requisição é direcionada ao servidor de borda geograficamente mais próximo, diminuindo a latência e o tempo de carregamento. Para otimizar esse processo e lidar com picos de audiência e a vasta quantidade de solicitações simultâneas, o balanceamento de carga torna-se indispensável. Ele distribui o tráfego de forma inteligente entre os diversos servidores de conteúdo de vídeo disponíveis, evitando que um único ponto da rede fique sobrecarregado. Essa distribuição equitativa é fundamental para prevenir gargalos, garantir a estabilidade do serviço e, consequentemente, otimizar a experiência de visualização para os consumidores.

Sendo assim, o objetivo principal deste trabalho é elaborar, implementar e analisar um algoritmo de balanceamento de carga para servidores de distribuição de conteúdo com base em monitoramento de memória, comparando-o com estratégias convencionais como \textit{Round-Robin} e Seleção Aleatória, focando na otimização do desempenho no contexto do provedor de serviço.

\section{O problema} \label{sec:problema}

Em sistemas distribuídos modernos caracterizados pela replicação de serviços, o balanceador de carga é peça fundamental no processo de maximização do aproveitamento de recursos computacionais e entende-se que sua efetividade está na capacidade de não sobrecarregar nenhum serviço mesmo em cenários de alta demanda \cite{dos2019balanceamento}. De modo simplificado, seu funcionamento consiste em, além de centralizar o roteamento de requisições, oferecer uma abstração para os consumidores, distribuir a carga entre vários servidores de modo a evitar sobrecarga em um ou mais deles. Nesse aspecto, podem ser empregadas várias estratégias para realizar esse roteamento: Seleção aleatória, \textit{Round-Robin}, seleção de um servidor baseada em algoritmos de eleição desenvolvidos para sistemas distribuídos, dentre outros. Enquanto são as mais utilizadas, a seleção aleatória e \textit{Round-Robin} são heurísticas simples de serem implementadas, mas que podem eventualmente causar sobrecarga em um ou mais servidores. A abordagem que contorna esse problema é justamente a eleição do servidor mais apropriado para lidar com a requisição de acordo com parâmetros personalizados que evidenciam, usando alguma estratégia, qual é o servidor a ser escolhido. A escolha dessa estratégia é a área-chave da atuação da proposta deste trabalho.

Além disso, dispositivos de armazenamento persistente (também conhecidos como dispositivos de memória secundária) são muitas vezes pontos de gargalo em sistemas, sobretudo sistemas de distribuição de conteúdo em \textit{streaming}, que exigem acesso aleatório à memória \cite{sivathanu2013IOWorkloads}. Esses dispositivos de armazenamento persistente exibem custos arquiteturais altos \cite{arpaciDusseau98}. Nesse sentido, em \textit{Data Centers} de larga escala, falhas de dispositivos de memória secundária são a regra e não a exceção \cite{wuEtalIOWorkload}, o que faz com que esses sistemas de alta escalabilidade tenham que ser tolerantes a falhas dos dispositivos de memória secundária. Dessa forma, entende-se que a métrica de taxa de uso de memória secundária é de grande valia no contexto da tomada de decisão de balanceadores de carga de requisições de distribuição de conteúdo. Em decorrência disso, para escalar sistemas na \textit{Web} de forma sólida, uma técnica muito difundida por projetistas de sistemas é o armazenamento de dados na memória principal. A redundância passou a ser característica de sistemas escaláveis na \textit{Web}. Alguns fatores que contribuiíam para isso foram: A evolução do número de usuários finais; O consequente surgimento do \textit{Big Data} gerando paradigmas revolucionadores no contexto de Bancos de Dados como os bancos não-relacionais \cite{riseofnewSQL}; O uso em massa das CDNs; A popularização da técnica de virtualização como uma forma de obter confiabilidade nos sistemas; dentre outros. Todos esses fatores corroboraram para que a memória principal, sob o seu principal benefício (a velocidade de acesso), fosse utilizada como método primário de armazenamento de dados persistentes, mesmo com a sua volatilidade intrínseca.

Dessa forma, considerando que o papel dos servidores de distribuição de conteúdo é recuperar e entregar dados (não envolvendo tanto processamento quanto outros tipos de serviços dependentes de Unidade Central de Processamento (CPU)), é possível conceber que a métrica do uso de memória (principal e secundária) é extremamente correlacionada com o desempenho desses servidores - considerando o contexto de múltiplas instâncias de servidores que mantêm seus próprios meios de armazenamento interno - cenário muito presente em CDNs com os servidores de borda (\textit{Edge Servers} ou \textit{Cache Servers}) \cite{surveyCDNArchitectures2025}.

Diante do exposto, evidencia-se uma lacuna crítica nas estratégias de balanceamento de carga convencionais. Abordagens como \textit{Round-Robin} ou seleção aleatória, ao ignorarem o estado atual dos recursos de memória, podem inadvertidamente direcionar requisições para nós computacionais já sobrecarregados, que estejam sofrendo com alta latência de leitura de disco ou próximos da exaustão de memória principal. Essa falha de percepção não só degrada o desempenho, como também aumenta o risco de falhas em cascata. Portanto, o problema central a ser atacado neste trabalho é a otimização da escolha do servidor em balanceadores de carga para sistemas de distribuição de conteúdo, desenvolvendo uma estratégia que incorpore o monitoramento de métricas de memória (primária e secundária) como fator decisório principal. O objetivo é criar um mecanismo de balanceamento mais inteligente e preditivo, capaz de evitar gargalos de desempenho antes que eles impactem o usuário final, aumentando assim a eficiência e a resiliência do sistema.

\section{Justificativa}

A eficiência dos Sistemas de Distribuição de Conteúdo (CDNs) é intrinsecamente relacionada à capacidade de seus servidores em entregar dados armazenados estaticamente de forma rápida aos usuários finais. Nesse contexto, a memória - tanto a principal utilizada em \textit{cache} quanto em disco - emerge como um recurso crítico, cujo estado reflete diretamente a aptidão de um servidor para responder a novas requisições de conteúdo.

Portanto, a proposta da elaboração e implementação de um algoritmo de balanceamento de carga justifica-se pela necessidade premente de explorar uma estratégia mais especializada e diretamente alinhada com a natureza operacional dos servidores de distribuição de conteúdo. Uma abordagem que utilize o monitoramento sistemático e exclusivo de memória dos servidores como critério para distribuição de carga visa preencher uma lacuna significativa e acredita-se que tal abordagem pode resultar em um sistema mais eficiente, capaz de otimizar a utilização de recursos de \textit{cache}, prevenir a sobrecarga dos servidores de forma mais eficaz e, consequentemente, aprimorar a latência e a experiência do usuário final na entrega do conteúdo.

\section{Objetivos}

Objetivo geral: Elaborar, implementar e analisar um algoritmo de balanceamento de carga para servidores de distribuição de conteúdo com base em monitoramento de memória de servidor de distribuição de conteúdo, comparando-o com estratégias convencionais como \textit{Round-Robin} e Seleção Aleatória;

Como objetivos específicos pode-se listar os seguintes:

\begin{itemize}
    \item Estudar as principais arquiteturas usadas em servidores de \textit{cache} de CDNs;
    \item Estudar os principais modelos de persistência de dados em servidores de distribuição de conteúdo;
    \item Revisar métodos de monitoramento de memória secundária;
    \item Revisar métodos de monitoramento de memória principal;
    \item Listar requisitos topológicos da rede e da estrutura dos servidores que definirão o contexto da aplicação da metodologia;
    \item Levantar estratégias de transmissão contínua de dados de monitoramento de servidores de distribuição de conteúdo para o serviço de balanceador de carga;
    \item Estabelecer, conceitualmente, algoritmo de escolha de servidores (por parte do balanceador de carga) com base nos dados históricos fornecidos;
    \item Implementar uma aplicação simples de balanceador de carga, e junto com ela o algoritmo criado;
    \item Avaliar os resultados obtidos por meio de testes de carga, tendo como principais métricas a latência de resposta, número e duração de travamentos de vídeo, e qualidade de reprodução; 
\end{itemize}

\section{Hipóteses}

Com base nos objetivos propostos e na fundamentação teórica que será apresentada, este trabalho parte das seguintes hipóteses a serem validadas experimentalmente:

\textbf{Hipótese 1:} O algoritmo de balanceamento baseado em monitoramento de memória apresentará desempenho superior aos métodos \textit{Round-Robin} e Aleatório em cenários heterogêneos, onde os servidores possuem capacidades distintas de memória e processamento.

\textbf{Hipótese 2:} O intervalo de atualização de 6 segundos das métricas de monitoramento fornecerá melhores resultados que o intervalo de 30 segundos, permitindo que o balanceador responda mais rapidamente às mudanças de estado dos servidores.

\textbf{Hipótese 3:} O algoritmo proposto será robusto mesmo com informações moderadamente desatualizadas, operando eficientemente com intervalos de atualização maiores (30 segundos), demonstrando sua viabilidade em sistemas de larga escala onde atualizações muito frequentes podem ser impraticáveis.

Essas hipóteses serão testadas por meio de experimentos controlados que simularão diferentes cenários de carga e heterogeneidade de servidores. Tanto os experimentos quanto os conceitos relacionados (intervalo de atualização, monitoramento de memória, desatualização de informações, etc.) serão descritos nos capítulos subsequentes.

\section{Metodologia e Estrutura}

Este trabalho consiste na realização das três principais tarefas a seguir: Apresentação e contextualização do problema; Realização de uma revisão de literatura com os objetivos de fundamentar teoricamente o tema e fornecer uma compreensão sólida; Proposição de um modelo de implementação com todos os detalhes técnicos e conceituais necessários. O restante do trabalho é organizado da seguinte maneira: O capítulo 2 apresenta contextualização e fundamentação teórica necessária, o capítulo 3 define a proposta deste trabalho e o capítulo 4 descreve o experimento a ser realizado cujos resultados serão apresentados no capítulo 5. Por fim, o capítulo 6 fornece uma conclusão acerca do experimento e dos resultados.
